Une cycliste (avec son vélo, \qty{\v{m}}{kg}) monte une côte dont le sommet est \qty{\v{h}}{m} plus haut que le bas. La montée dure \qty{\v{t}}{s}. On utilise $g = \qty{9.8}{N/kg}$.

\begin{enumerate}[label=\alph*)]
    \item Quelle énergie potentielle gravitationnelle la cycliste gagne-t-elle ?
    \item Quelle puissance utile moyenne développe-t-elle pendant la montée ?
    \item Le rendement de son corps est de \qty{25}{\percent}. Quelle quantité d'énergie a-t-elle réellement dépensée ?
\end{enumerate}
